\section{Latar Belakang}
\TemplateTodo{Jelaskan konteks umum topik, kondisi ideal, masalah nyata, kesenjangan penelitian atau kebutuhan praktis, serta alasan penelitian perlu dilakukan.}

Pertumbuhan pengguna \emph{smartphone} di Indonesia meningkat pesat setiap tahun. Berdasarkan laporan We Are Social, jumlah pengguna \emph{smartphone} di Indonesia mencapai 212,9 juta jiwa pada tahun 2024 \cite{manning2008}. Peningkatan ini mendorong berkembangnya layanan berbasis aplikasi \emph{mobile}, termasuk layanan transportasi dan pesan-antar makanan \emph{online}. Aplikasi GoFood sebagai salah satu layanan terbesar di Indonesia memiliki jutaan ulasan di Google Play Store yang mencerminkan opini dan pengalaman pengguna.

Ulasan pengguna mengandung informasi berharga bagi pengembang aplikasi untuk meningkatkan kualitas layanan. Namun, menganalisis ribuan ulasan secara manual membutuhkan waktu dan tenaga yang tidak sedikit. Oleh karena itu, diperlukan pendekatan otomatis untuk mengklasifikasikan sentimen ulasan ke dalam kategori positif dan negatif.

Analisis sentimen merupakan salah satu cabang \emph{Natural Language Processing} (NLP) yang dapat digunakan untuk mengekstrak opini dari teks secara otomatis \cite{feldman2013}. Metode Na\"ive Bayes Classifier (NBC) dipilih karena sederhana, cepat, dan telah terbukti efektif dalam klasifikasi teks \cite{go2019}. Penelitian ini menerapkan NBC dengan fitur TF-IDF untuk mengklasifikasikan sentimen ulasan aplikasi GoFood.
